\documentclass[11pt]{article}

\usepackage[utf8]{inputenc}
\usepackage[T1]{fontenc}
\usepackage{lmodern}
\usepackage{geometry}
\usepackage{amsmath,amssymb,amsfonts,amsthm,mathtools}
\usepackage{bm}
\usepackage{enumitem}
\usepackage{microtype}
\usepackage{hyperref}
\usepackage{titlesec}
\usepackage{booktabs}
\usepackage{array}
\usepackage{setspace}
\usepackage{mathrsfs}
\usepackage{physics}

\geometry{margin=1in}
\hypersetup{colorlinks=true, linkcolor=black, urlcolor=blue, citecolor=black}
\setlist[enumerate]{topsep=0.4em,itemsep=0.35em}
\setlist[itemize]{topsep=0.3em,itemsep=0.25em}
\titleformat{\section}{\large\bfseries}{\thesection.}{0.5em}{}
\titleformat{\subsection}{\normalsize\bfseries}{\thesubsection.}{0.5em}{}
\setstretch{1.06}

\newcommand{\Aether}{\AE{}ther}
\newcommand{\Aflow}{\AE{}ther-flow}
\newcommand{\TheoryName}{\Aflow{} Interpretation of Relativity}
\newcommand{\TheoryProgram}{\Aether{} / \Aflow{} framework}
\newcommand{\CompletedMetric}{g^{\sharp}}
\newcommand{\SelectorCore}{\mathfrak C^{\mathrm{sel,split}}}

\newcommand{\Car}{\mathrm{car}}
\newcommand{\Cur}{\mathrm{cur}}
\newcommand{\Hom}{\mathrm{hom}}

\newcommand{\ForSpace}[1]{\mathcal I_{#1}^{\mathrm{for}}}
\newcommand{\PrimShell}{\mathcal S_{\mathrm{prim}}}
\newcommand{\Reduction}[1]{\mathcal R_{#1}}
\newcommand{\CurrentCarrier}[1]{\mathfrak C_{#1}^{\mathrm{cur}}}
\newcommand{\EqCarrier}[1]{\mathfrak C_{#1}^{\mathrm{eq}}}
\newcommand{\MetricOut}{\mathscr G_{\mathrm{out}}}
\newcommand{\MatterOut}{\mathscr T_{\mathrm{out}}}

\newcommand{\VisibleAffine}[1]{\widehat X_{#1}}
\newcommand{\DataSpace}[1]{\mathcal Y_{#1}^{\mathrm{obs}}}
\newcommand{\AmbientTarget}[1]{E_{#1}^{\mathrm{amb}}}

\newcommand{\GenSpace}[1]{\mathcal G_{#1}^{\mathrm{sel}}}
\newcommand{\GenBody}[1]{\mathcal B_{#1}^{\mathrm{sel}}}
\newcommand{\GenData}[1]{\Lambda_{#1}^{\mathrm{sel,dat}}}
\newcommand{\GenShell}[1]{\Lambda_{#1}^{\mathrm{sel,sh}}}
\newcommand{\HiddenSel}[1]{\Sigma_{#1}^{\mathrm{hid}}}
\newcommand{\ZeroHiddenSlice}[1]{\mathcal Z_{#1}^{0}}
\newcommand{\ZeroHiddenVec}[1]{V_{#1}^{0}}
\newcommand{\GapSpace}[1]{H_{#1}}
\newcommand{\GapBody}[1]{D_{#1}}
\newcommand{\HiddenOperator}[1]{K_{#1}^{\mathrm{hid}}}
\newcommand{\GenSym}[1]{\mathbb S_{#1}}
\newcommand{\GenTime}[1]{\vartheta_{#1}^{\mathrm{sel}}}
\newcommand{\GenEmbed}[1]{\chi_{#1}}
\newcommand{\Kernel}[1]{K_{#1}^{0}}
\newcommand{\StateComp}[1]{P_{#1}^{\mathrm{co}}}

\newcommand{\SelCoreSpace}[1]{\mathcal U_{#1}^{\mathrm{sel,co}}}
\newcommand{\SelCoreBody}[1]{\mathcal B_{#1}^{\mathrm{sel,co}}}
\newcommand{\SelCoreData}[1]{\Lambda_{#1}^{\mathrm{sel,co,dat}}}
\newcommand{\SelCoreShell}[1]{\Lambda_{#1}^{\mathrm{sel,co,sh}}}
\newcommand{\SelHiddenData}[1]{\Gamma_{#1}^{\mathrm{sel,hid,dat}}}
\newcommand{\SelHiddenShell}[1]{\Gamma_{#1}^{\mathrm{sel,hid,sh}}}
\newcommand{\SelCoreSym}[1]{\mathbb T_{#1}}
\newcommand{\SelCoreTime}[1]{\vartheta_{#1}^{\mathrm{sel,co}}}
\newcommand{\SelCoreEmbed}[1]{\xi_{#1}}
\newcommand{\SelCoreProj}[1]{P_{#1}^{\mathrm{sel,co}}}

\newtheorem{definition}{Definition}
\newtheorem{lemma}{Lemma}
\newtheorem{proposition}{Proposition}
\newtheorem{remark}{Remark}
\newtheorem{corollary}{Corollary}
\newtheorem{theorem}{Theorem}

\title{\textbf{The \TheoryName{}}\\[0.4em]
\large Primitive-Reservoir Observer-Reduction\\
\large Zero-Hidden Selector Generator Law\\
\large Collapse to a Split Zero-Hidden Selector-Shell Core}
\author{}
\date{}

\begin{document}

\maketitle

\begin{abstract}
The current benchmark-facing theorem on the active primitive-reservoir observer-reduction line already removed the combined response substrate law as the live unexplained stopping point: that law is now a canonical and necessary reduction of a deeper zero-hidden selector generator law. The remaining honest burden therefore sharpened to the selector law itself. The active routing safeguard then fixed two admissible routes below that frontier: either force the selector law from still more primitive explicit substrate structure, or prove a genuinely smaller theorem-level collapse strictly inside that selector law.

The present note takes the second route. It proves that the full zero-hidden selector generator law carries no extra benchmark-facing freedom beyond a smaller split zero-hidden selector-shell core \(\SelectorCore\). For each sector that smaller datum keeps only the zero-hidden state space and body, the affine observer-data and shell-response readouts on that zero-hidden sector, the hidden-to-data and hidden-to-shell coupling maps, the hidden body with positive hidden-sector operator, the induced symmetry and time-reversal structure, the exact zero-response shell realization, fixed-data shell solvability, the inert zero-hidden kernel and projector, and the split coercive control on data-invisible zero-hidden and hidden variations. What it omits as independent primitive data are the larger ambient selector-state space and the self-adjoint idempotent hidden selector itself. Those are shown to be canonically induced once the split core is fixed.

This is a genuine smaller theorem-level collapse strictly inside the current selector frontier. The full selector-state presentation is recovered canonically as the orthogonal product of the zero-hidden and hidden sectors, the hidden selector is recovered as the product projection, and the full affine selector readouts are recovered from the zero-hidden readouts together with the hidden coupling maps. Consequently all current downstream uses of the selector law, including the already recorded forcing of the combined response substrate law, factor through \(\SelectorCore\). The benchmark boundary remains fixed throughout: the one operative metric is still exactly \(\CompletedMetric\), the ordinary matter route is unchanged, there is no second operative metric, the low-energy matter law is not enlarged, and stronger first-principles promotion language remains paused.
\end{abstract}

\section{Introduction}

The active derivational program of \TheoryName{} has reached another stage at which depth alone does not count as progress. The primitive-observer-reduction datum theorem, the defect-fiber factorization theorem, the one-metric output theorem, the chain-forcing theorem, the selector theorem, the stationary-construction theorem, the explicit-package theorem, the primitive-normal-form theorem, the exact-shell-affine-nucleus collapse theorem, the completion-core trace-collapse theorem, the ambient-dynamics forcing theorem, and the current selector-law forcing theorem have already discharged what was honestly open at their respective layers. Once the combined response substrate law is itself forced by the selector law, another same-output combined-response relay, ambient-package relay, or selector-law relay that leaves the selector law unchanged no longer counts as the honest next benchmark-facing move.

The routing layer therefore permitted two disciplined routes below the selector frontier. One could derive or uniquely force the full zero-hidden selector generator law from still more primitive explicit substrate structure in a way that genuinely changes benchmark-facing status. Or one could show that the current selector law still carries redundant presentation freedom and collapse it to a smaller benchmark-facing datum strictly inside that same law. The present note takes the second route.

Its gain is narrower than a completed first-principles substrate derivation and stronger than another deeper same-output relay. The key observation is that the full selector law already singles out its own split sectors: the zero-hidden sector, the hidden sector, and the affine readout couplings between them. Once those split pieces are fixed together with the exact shell realization, solvability statement, inert-kernel projector, and coercive control, the larger ambient selector-state presentation and the hidden selector idempotent cease to be independent benchmark-facing data. They are reconstructed canonically as product presentation.

\paragraph{Framework and claim boundary.}
\TheoryProgram{} denotes the broader \Aether{} / \Aflow{} framework built on the claims that the \Aether{} is the underlying four-dimensional substrate of reality, the \Aflow{} is its intrinsic ordered motion, observed three-dimensional space is the local experiential slice of that deeper substrate, S-time is the experienced order of change arising from matter, light, and the \Aflow{}, and observed expansion is the three-dimensional appearance of deeper four-dimensional ordered motion. Within that framework, \TheoryName{} denotes the exact-closure theory in which the effective gravitational dynamics are adopted to be exactly Einsteinian with universal matter coupling. In the active sequence, \emph{adoption} means use of that established relativistic dynamics without claiming substrate derivation, while \emph{derivation} is reserved for a first-principles recovery from explicit substrate variables.

The present note stays strictly inside that boundary. It does not alter the benchmark exact-closure package. It does not introduce a second operative metric or a new low-energy matter law. It only proves that the current selector-law burden collapses further to a smaller split selector-shell core at benchmark-facing scope.

\section{Fixed Primitive Continuation Data}

\begin{definition}[Recorded primitive continuation data]
\label{def:fixed}
Fix the following continuation data from the active primitive-reservoir observer-reduction line.
\begin{enumerate}
    \item For each sector label \(\sigma\in\{\Car,\Cur,\Hom\}\), a primitive forcing shell
    \[
    \ForSpace{\sigma},
    \]
    a visible reduction map
    \[
    \Reduction{\sigma}:\ForSpace{\sigma}\to X_{\sigma},
    \]
    and shell-level current and equilibrium carriers
    \[
    \CurrentCarrier{\sigma}:\ForSpace{\sigma}\to Z_{\sigma}^{\mathrm{cur}},
    \qquad
    \EqCarrier{\sigma}:\ForSpace{\sigma}\to Z_{\sigma}^{\mathrm{eq}}.
    \]
    \item The product shell
    \[
    \PrimShell
    \equiv
    \ForSpace{\Car}\times \ForSpace{\Cur}\times \ForSpace{\Hom}.
    \]
    \item The recorded metric and ordinary-matter outputs
    \[
    \MetricOut:\PrimShell\to\{\text{observer metrics}\},
    \qquad
    \MatterOut:\PrimShell\times\Psi\to\{\text{observer stress tensors}\},
    \]
    satisfying
    \begin{equation}
    \MetricOut(\mathbf w)=\CompletedMetric,
    \qquad
    \MatterOut(\mathbf w,\psi)
    =
    T_{\mu\nu}^{\mathrm{matter}}[\CompletedMetric,\psi]
    \label{eq:fixedoutputs}
    \end{equation}
    for every \(\mathbf w\in\PrimShell\) and every matter configuration \(\psi\in\Psi\).
\end{enumerate}
\end{definition}

\begin{remark}
Definition \ref{def:fixed} is not reopened here. The primitive shell data, the one operative metric, and the ordinary matter route remain fixed continuation data. The present note addresses only the benchmark-facing internal compression of the current selector-law frontier.
\end{remark}

\section{Recorded Full Selector Law}

\begin{definition}[Zero-hidden selector generator law]
\label{def:generator}
Fix \(\sigma\in\{\Car,\Cur,\Hom\}\). A \emph{zero-hidden selector generator law} for sector \(\sigma\) consists of the following data.
\begin{enumerate}
    \item A finite-dimensional Euclidean selector-state space
    \[
    \GenSpace{\sigma},
    \]
    the observer-data space
    \[
    \DataSpace{\sigma}
    \equiv
    \VisibleAffine{\sigma}\times Z_{\sigma}^{\mathrm{cur}}\times Z_{\sigma}^{\mathrm{eq}},
    \]
    an ambient shell-response space
    \[
    \AmbientTarget{\sigma},
    \]
    and a compact convex selector-state body
    \[
    \GenBody{\sigma}\subset\GenSpace{\sigma}
    \]
    with nonempty interior.
    \item An affine observer-data readout
    \[
    \GenData{\sigma}:\GenSpace{\sigma}\to\DataSpace{\sigma}
    \]
    and an affine shell-response readout
    \[
    \GenShell{\sigma}:\GenSpace{\sigma}\to\AmbientTarget{\sigma}.
    \]
    \item A self-adjoint idempotent hidden selector
    \[
    \HiddenSel{\sigma}:\GenSpace{\sigma}\to\GenSpace{\sigma},
    \qquad
    {\HiddenSel{\sigma}}^{2}=\HiddenSel{\sigma},
    \qquad
    {\HiddenSel{\sigma}}^{*}=\HiddenSel{\sigma}.
    \]
    Define the selector-zero locus
    \[
    \ZeroHiddenSlice{\sigma}
    \equiv
    \ker\HiddenSel{\sigma},
    \]
    its translation space
    \[
    \ZeroHiddenVec{\sigma}
    \equiv
    \ker\HiddenSel{\sigma},
    \]
    and the hidden selector image
    \[
    \GapSpace{\sigma}
    \equiv
    \operatorname{im}\HiddenSel{\sigma}.
    \]
    \item A compact convex hidden body
    \[
    \GapBody{\sigma}\subset\GapSpace{\sigma}
    \]
    with
    \[
    0\in\operatorname{int}\GapBody{\sigma},
    \]
    such that the selector-state body splits exactly as
    \begin{equation}
    \GenBody{\sigma}
    =
    \left\{
    z+\eta:
    z\in\GenBody{\sigma}\cap\ZeroHiddenSlice{\sigma},
    \quad
    \eta\in\GapBody{\sigma}
    \right\}.
    \label{eq:gensplitting}
    \end{equation}
    \item A positive self-adjoint hidden-sector operator
    \[
    \HiddenOperator{\sigma}:\GapSpace{\sigma}\to\GapSpace{\sigma}
    \]
    satisfying
    \begin{equation}
    \ev{\eta,\HiddenOperator{\sigma}\eta}
    \ge
    \kappa_{\sigma}^{\mathrm{hid}}\,
    \|\eta\|^{2}
    \qquad
    \text{for all }\eta\in\GapSpace{\sigma}
    \label{eq:gengap}
    \end{equation}
    for some \(\kappa_{\sigma}^{\mathrm{hid}}>0\).
    \item A compact symmetry group
    \[
    \GenSym{\sigma}\subset
    O(\GenSpace{\sigma})\times
    O(\DataSpace{\sigma})\times
    O(\AmbientTarget{\sigma})
    \]
    and an orthogonal involution
    \[
    \GenTime{\sigma}\in
    O(\GenSpace{\sigma})\times
    O(\DataSpace{\sigma})\times
    O(\AmbientTarget{\sigma})
    \]
    preserving \(\GenBody{\sigma}\), \(\GenData{\sigma}\), \(\GenShell{\sigma}\), \(\HiddenSel{\sigma}\), \(\GapBody{\sigma}\), and \(\HiddenOperator{\sigma}\).
    \item A smooth embedding
    \[
    \jmath_{\sigma}:X_{\sigma}\hookrightarrow\VisibleAffine{\sigma}
    \]
    and a smooth zero-response shell realization
    \[
    \GenEmbed{\sigma}:\ForSpace{\sigma}\hookrightarrow\GenSpace{\sigma}
    \]
    whose image is exactly
    \[
    \GenBody{\sigma}\cap
    \ZeroHiddenSlice{\sigma}\cap
    \bigl(\GenShell{\sigma}\bigr)^{-1}(0),
    \]
    and such that
    \begin{equation}
    \GenData{\sigma}\bigl(\GenEmbed{\sigma}(w)\bigr)
    =
    \bigl(
    \jmath_{\sigma}(\Reduction{\sigma}(w)),
    \CurrentCarrier{\sigma}(w),
    \EqCarrier{\sigma}(w)
    \bigr)
    \label{eq:gencompat}
    \end{equation}
    for every \(w\in\ForSpace{\sigma}\).
    \item Fixed-data shell solvability on the selector-zero locus:
    \begin{equation}
    \GenData{\sigma}\bigl(
    \GenBody{\sigma}\cap\ZeroHiddenSlice{\sigma}
    \bigr)
    =
    \GenData{\sigma}\Bigl(
    \GenBody{\sigma}\cap
    \ZeroHiddenSlice{\sigma}\cap
    \bigl(\GenShell{\sigma}\bigr)^{-1}(0)
    \Bigr).
    \label{eq:gensolvability}
    \end{equation}
    \item The inert zero-hidden kernel
    \[
    \Kernel{\sigma}
    \equiv
    \ker\bigl(D\GenData{\sigma}\big|_{\ZeroHiddenVec{\sigma}}\bigr)
    \cap
    \ker\bigl(D\GenShell{\sigma}\big|_{\ZeroHiddenVec{\sigma}}\bigr)
    \]
    and the orthogonal projector
    \[
    \StateComp{\sigma}:
    \ZeroHiddenVec{\sigma}\to
    {\Kernel{\sigma}}^{\perp}\cap\ZeroHiddenVec{\sigma}.
    \]
    \item Selector-law coercivity on data-invisible directions: there exists \(\kappa_{\sigma}^{\mathrm{sel}}>0\) such that for every
    \[
    w\in\ForSpace{\sigma},
    \qquad
    \delta g\in T_{\GenEmbed{\sigma}(w)}\GenSpace{\sigma}
    \]
    satisfying
    \[
    D\GenData{\sigma}\,\delta g=0,
    \]
    one has
    \begin{equation}
    \|D\GenShell{\sigma}\,\delta g\|^{2}
    +
    \ev{
    \HiddenSel{\sigma}\delta g,
    \HiddenOperator{\sigma}\HiddenSel{\sigma}\delta g
    }
    \ge
    \kappa_{\sigma}^{\mathrm{sel}}
    \left(
    \|\StateComp{\sigma}( (I-\HiddenSel{\sigma})\delta g )\|^{2}
    +
    \|\HiddenSel{\sigma}\delta g\|^{2}
    \right).
    \label{eq:gencoercive}
    \end{equation}
\end{enumerate}
\end{definition}

\begin{remark}
Definition \ref{def:generator} records the full selector-law frontier currently on the active line. The purpose of the present note is not to restate that law as another stopping point, but to identify which constituents there remain independently benchmark-facing once its split sectors are made explicit.
\end{remark}

\section{Split Zero-Hidden Selector-Shell Core}

\begin{definition}[Split zero-hidden selector-shell core with hidden gap]
\label{def:splitcore}
Fix \(\sigma\in\{\Car,\Cur,\Hom\}\). A \emph{split zero-hidden selector-shell core with hidden gap} for sector \(\sigma\) consists of the following data.
\begin{enumerate}
    \item A finite-dimensional Euclidean zero-hidden state space
    \[
    \SelCoreSpace{\sigma},
    \]
    the observer-data space \(\DataSpace{\sigma}\), the ambient shell-response space \(\AmbientTarget{\sigma}\), the hidden space \(\GapSpace{\sigma}\), and a compact convex zero-hidden body
    \[
    \SelCoreBody{\sigma}\subset\SelCoreSpace{\sigma}
    \]
    with nonempty interior.
    \item An affine zero-hidden observer-data readout
    \[
    \SelCoreData{\sigma}:\SelCoreSpace{\sigma}\to\DataSpace{\sigma}
    \]
    and an affine zero-hidden shell-response readout
    \[
    \SelCoreShell{\sigma}:\SelCoreSpace{\sigma}\to\AmbientTarget{\sigma}.
    \]
    \item Linear hidden-to-data and hidden-to-shell coupling maps
    \[
    \SelHiddenData{\sigma}:\GapSpace{\sigma}\to\DataSpace{\sigma},
    \qquad
    \SelHiddenShell{\sigma}:\GapSpace{\sigma}\to\AmbientTarget{\sigma}.
    \]
    \item A compact convex hidden body
    \[
    \GapBody{\sigma}\subset\GapSpace{\sigma}
    \]
    with
    \[
    0\in\operatorname{int}\GapBody{\sigma},
    \]
    together with a positive self-adjoint hidden-sector operator
    \[
    \HiddenOperator{\sigma}:\GapSpace{\sigma}\to\GapSpace{\sigma}
    \]
    satisfying \eqref{eq:gengap}.
    \item A compact symmetry group
    \[
    \SelCoreSym{\sigma}\subset
    O(\SelCoreSpace{\sigma})\times
    O(\DataSpace{\sigma})\times
    O(\AmbientTarget{\sigma})\times
    O(\GapSpace{\sigma})
    \]
    and an orthogonal involution
    \[
    \SelCoreTime{\sigma}\in
    O(\SelCoreSpace{\sigma})\times
    O(\DataSpace{\sigma})\times
    O(\AmbientTarget{\sigma})\times
    O(\GapSpace{\sigma})
    \]
    preserving \(\SelCoreBody{\sigma}\), \(\SelCoreData{\sigma}\), \(\SelCoreShell{\sigma}\), \(\SelHiddenData{\sigma}\), \(\SelHiddenShell{\sigma}\), \(\GapBody{\sigma}\), and \(\HiddenOperator{\sigma}\).
    \item A smooth embedding
    \[
    \jmath_{\sigma}:X_{\sigma}\hookrightarrow\VisibleAffine{\sigma}
    \]
    and a smooth zero-response shell realization
    \[
    \SelCoreEmbed{\sigma}:\ForSpace{\sigma}\hookrightarrow\SelCoreSpace{\sigma}
    \]
    whose image is exactly
    \[
    \SelCoreBody{\sigma}\cap
    \bigl(\SelCoreShell{\sigma}\bigr)^{-1}(0),
    \]
    and such that
    \begin{equation}
    \SelCoreData{\sigma}\bigl(\SelCoreEmbed{\sigma}(w)\bigr)
    =
    \bigl(
    \jmath_{\sigma}(\Reduction{\sigma}(w)),
    \CurrentCarrier{\sigma}(w),
    \EqCarrier{\sigma}(w)
    \bigr)
    \label{eq:splitcompat}
    \end{equation}
    for every \(w\in\ForSpace{\sigma}\).
    \item Fixed-data shell solvability on the zero-hidden sector:
    \begin{equation}
    \SelCoreData{\sigma}\bigl(\SelCoreBody{\sigma}\bigr)
    =
    \SelCoreData{\sigma}\Bigl(
    \SelCoreBody{\sigma}\cap
    \bigl(\SelCoreShell{\sigma}\bigr)^{-1}(0)
    \Bigr).
    \label{eq:splitsolvability}
    \end{equation}
    \item The inert zero-hidden kernel
    \[
    \Kernel{\sigma}
    \equiv
    \ker\bigl(D\SelCoreData{\sigma}\bigr)
    \cap
    \ker\bigl(D\SelCoreShell{\sigma}\bigr)
    \]
    and the orthogonal projector
    \[
    \SelCoreProj{\sigma}:
    \SelCoreSpace{\sigma}\to
    {\Kernel{\sigma}}^{\perp}\cap\SelCoreSpace{\sigma}.
    \]
    \item Split-core coercivity on data-invisible zero-hidden and hidden variations: there exists \(\kappa_{\sigma}^{\mathrm{split}}>0\) such that for every
    \[
    w\in\ForSpace{\sigma},
    \qquad
    \delta u\in T_{\SelCoreEmbed{\sigma}(w)}\SelCoreSpace{\sigma},
    \qquad
    \eta\in\GapSpace{\sigma}
    \]
    satisfying
    \[
    D\SelCoreData{\sigma}\,\delta u+\SelHiddenData{\sigma}\eta=0,
    \]
    one has
    \begin{equation}
    \bigl\|
    D\SelCoreShell{\sigma}\,\delta u+\SelHiddenShell{\sigma}\eta
    \bigr\|^{2}
    +
    \ev{\eta,\HiddenOperator{\sigma}\eta}
    \ge
    \kappa_{\sigma}^{\mathrm{split}}
    \left(
    \|\SelCoreProj{\sigma}\delta u\|^{2}
    +
    \|\eta\|^{2}
    \right).
    \label{eq:splitcoercive}
    \end{equation}
\end{enumerate}
\end{definition}

\begin{remark}
Definition \ref{def:splitcore} is strictly smaller than Definition \ref{def:generator}. It keeps the benchmark-relevant split sectors and their couplings, but it does not treat the larger ambient selector-state space or the hidden selector projector as independent primitive data.
\end{remark}

\section{Canonical Split Core Induced by the Selector Law}

\begin{definition}[Canonical split core induced by the selector law]
\label{def:canonical}
Assume Definition \ref{def:generator}. Define the canonical split zero-hidden selector-shell core by:
\begin{enumerate}
    \item
    \[
    \SelCoreSpace{\sigma}\equiv\ZeroHiddenVec{\sigma},
    \qquad
    \SelCoreBody{\sigma}\equiv
    \GenBody{\sigma}\cap\ZeroHiddenSlice{\sigma}.
    \]
    \item
    \[
    \SelCoreData{\sigma}
    \equiv
    \GenData{\sigma}\big|_{\ZeroHiddenVec{\sigma}},
    \qquad
    \SelCoreShell{\sigma}
    \equiv
    \GenShell{\sigma}\big|_{\ZeroHiddenVec{\sigma}}.
    \]
    \item The hidden coupling maps are
    \begin{equation}
    \SelHiddenData{\sigma}(\eta)
    \equiv
    \GenData{\sigma}(\eta)-\GenData{\sigma}(0),
    \qquad
    \SelHiddenShell{\sigma}(\eta)
    \equiv
    \GenShell{\sigma}(\eta)-\GenShell{\sigma}(0)
    \label{eq:hiddencouplings}
    \end{equation}
    for \(\eta\in\GapSpace{\sigma}\).
    \item The hidden body and hidden operator are the same \(\GapBody{\sigma}\) and \(\HiddenOperator{\sigma}\) from Definition \ref{def:generator}.
    \item The symmetry and time-reversal data are the restricted selector symmetries:
    \[
    \SelCoreSym{\sigma}
    \equiv
    \left\{
    \bigl(
    g_{0},g_{\mathrm{dat}},g_{\mathrm{sh}},g_{\mathrm{hid}}
    \bigr):
    (g_{\mathrm{sel}},g_{\mathrm{dat}},g_{\mathrm{sh}})
    \in\GenSym{\sigma}
    \right\},
    \]
    where
    \[
    g_{0}\equiv g_{\mathrm{sel}}\big|_{\ZeroHiddenVec{\sigma}},
    \qquad
    g_{\mathrm{hid}}\equiv g_{\mathrm{sel}}\big|_{\GapSpace{\sigma}},
    \]
    and
    \[
    \SelCoreTime{\sigma}
    \equiv
    \bigl(
    \GenTime{\sigma}\big|_{\ZeroHiddenVec{\sigma}},
    \GenTime{\sigma}\big|_{\DataSpace{\sigma}},
    \GenTime{\sigma}\big|_{\AmbientTarget{\sigma}},
    \GenTime{\sigma}\big|_{\GapSpace{\sigma}}
    \bigr).
    \]
    \item The exact shell realization is
    \[
    \SelCoreEmbed{\sigma}\equiv\GenEmbed{\sigma}.
    \]
    \item The inert kernel and projector are the same:
    \[
    \Kernel{\sigma}
    \equiv
    \ker\bigl(D\SelCoreData{\sigma}\bigr)
    \cap
    \ker\bigl(D\SelCoreShell{\sigma}\bigr),
    \qquad
    \SelCoreProj{\sigma}\equiv\StateComp{\sigma}.
    \]
\end{enumerate}
\end{definition}

\begin{proposition}[The selector law induces a split selector-shell core]
\label{prop:induced}
Assume Definitions \ref{def:generator}, \ref{def:splitcore}, and \ref{def:canonical}. Then the canonical data of Definition \ref{def:canonical} satisfy every requirement of Definition \ref{def:splitcore}.
\end{proposition}

\begin{proof}
Because \(\HiddenSel{\sigma}\) is self-adjoint and idempotent, the selector-state space splits orthogonally as
\[
\GenSpace{\sigma}
=
\ZeroHiddenVec{\sigma}\oplus\GapSpace{\sigma}.
\]
Hence \(\SelCoreSpace{\sigma}=\ZeroHiddenVec{\sigma}\) is a Euclidean space and \(\SelCoreBody{\sigma}=\GenBody{\sigma}\cap\ZeroHiddenSlice{\sigma}\) is compact and convex. By the exact splitting \eqref{eq:gensplitting} with \(0\in\operatorname{int}\GapBody{\sigma}\), the zero-hidden slice carries nonempty interior inside its own translation space, so \(\SelCoreBody{\sigma}\) has nonempty interior in \(\SelCoreSpace{\sigma}\).

The restrictions \(\SelCoreData{\sigma}\) and \(\SelCoreShell{\sigma}\) are affine because \(\GenData{\sigma}\) and \(\GenShell{\sigma}\) are affine on \(\GenSpace{\sigma}\). The coupling maps \(\SelHiddenData{\sigma}\) and \(\SelHiddenShell{\sigma}\) are linear because they are obtained by subtracting the constant affine offsets in \eqref{eq:hiddencouplings}. The hidden body and hidden operator already satisfy the required positivity and interior conditions from Definition \ref{def:generator}.

Since the selector symmetries preserve \(\HiddenSel{\sigma}\), they preserve both \(\ker\HiddenSel{\sigma}\) and \(\operatorname{im}\HiddenSel{\sigma}\). Therefore the restricted symmetry and time-reversal data are well defined and preserve the split-core constituents. The shell realization statement in Definition \ref{def:splitcore} is immediate because \(\GenEmbed{\sigma}(\ForSpace{\sigma})\) is exactly
\[
\GenBody{\sigma}\cap\ZeroHiddenSlice{\sigma}\cap
\bigl(\GenShell{\sigma}\bigr)^{-1}(0)
=
\SelCoreBody{\sigma}\cap
\bigl(\SelCoreShell{\sigma}\bigr)^{-1}(0),
\]
and \eqref{eq:splitcompat} is just \eqref{eq:gencompat}.

Fixed-data shell solvability \eqref{eq:splitsolvability} is precisely the selector-law solvability statement \eqref{eq:gensolvability} restricted to the zero-hidden sector. The inert kernel and projector are the same by construction.

For coercivity, let \(w\in\ForSpace{\sigma}\), \(\delta u\in T_{\SelCoreEmbed{\sigma}(w)}\SelCoreSpace{\sigma}\), and \(\eta\in\GapSpace{\sigma}\) satisfy
\[
D\SelCoreData{\sigma}\,\delta u+\SelHiddenData{\sigma}\eta=0.
\]
Set \(\delta g\equiv \delta u+\eta\in\GenSpace{\sigma}\). Then
\[
D\GenData{\sigma}\,\delta g
=
D\SelCoreData{\sigma}\,\delta u+\SelHiddenData{\sigma}\eta
=
0,
\]
while
\[
D\GenShell{\sigma}\,\delta g
=
D\SelCoreShell{\sigma}\,\delta u+\SelHiddenShell{\sigma}\eta,
\qquad
\HiddenSel{\sigma}\delta g=\eta,
\qquad
(I-\HiddenSel{\sigma})\delta g=\delta u.
\]
Applying \eqref{eq:gencoercive} therefore gives exactly \eqref{eq:splitcoercive}. All requirements of Definition \ref{def:splitcore} are satisfied.
\end{proof}

\section{Reconstructing the Full Selector Law from the Split Core}

\begin{definition}[Canonical selector law recovered from a split core]
\label{def:recovered}
Assume Definition \ref{def:splitcore}. Define the recovered selector law as follows.
\begin{enumerate}
    \item The selector-state space is the orthogonal product
    \begin{equation}
    \GenSpace{\sigma}
    \equiv
    \SelCoreSpace{\sigma}\oplus\GapSpace{\sigma}.
    \label{eq:recspace}
    \end{equation}
    Write elements as \(u+\eta\) with \(u\in\SelCoreSpace{\sigma}\) and \(\eta\in\GapSpace{\sigma}\).
    \item The selector-state body is
    \begin{equation}
    \GenBody{\sigma}
    \equiv
    \left\{
    u+\eta:
    u\in\SelCoreBody{\sigma},
    \ \eta\in\GapBody{\sigma}
    \right\}.
    \label{eq:recbody}
    \end{equation}
    \item The affine selector readouts are
    \begin{equation}
    \GenData{\sigma}(u+\eta)
    \equiv
    \SelCoreData{\sigma}(u)+\SelHiddenData{\sigma}(\eta),
    \qquad
    \GenShell{\sigma}(u+\eta)
    \equiv
    \SelCoreShell{\sigma}(u)+\SelHiddenShell{\sigma}(\eta).
    \label{eq:recreadouts}
    \end{equation}
    \item The hidden selector is the orthogonal product projection
    \begin{equation}
    \HiddenSel{\sigma}(u+\eta)\equiv \eta.
    \label{eq:reselector}
    \end{equation}
    Consequently
    \[
    \ZeroHiddenSlice{\sigma}
    \equiv
    \SelCoreSpace{\sigma},
    \qquad
    \ZeroHiddenVec{\sigma}
    \equiv
    \SelCoreSpace{\sigma},
    \qquad
    \operatorname{im}\HiddenSel{\sigma}
    \equiv
    \GapSpace{\sigma}.
    \]
    \item The symmetry and time-reversal data are
    \[
    \GenSym{\sigma}
    \equiv
    \left\{
    \bigl(
    g_{0}\oplus g_{\mathrm{hid}},
    g_{\mathrm{dat}},
    g_{\mathrm{sh}}
    \bigr):
    \bigl(
    g_{0},g_{\mathrm{dat}},g_{\mathrm{sh}},g_{\mathrm{hid}}
    \bigr)\in\SelCoreSym{\sigma}
    \right\},
    \]
    and
    \[
    \GenTime{\sigma}
    \equiv
    \bigl(
    \SelCoreTime{\sigma}\big|_{\SelCoreSpace{\sigma}}
    \oplus
    \SelCoreTime{\sigma}\big|_{\GapSpace{\sigma}},
    \SelCoreTime{\sigma}\big|_{\DataSpace{\sigma}},
    \SelCoreTime{\sigma}\big|_{\AmbientTarget{\sigma}}
    \bigr).
    \]
    \item The exact shell realization is
    \begin{equation}
    \GenEmbed{\sigma}(w)\equiv \SelCoreEmbed{\sigma}(w)+0.
    \label{eq:recembed}
    \end{equation}
    \item The inert zero-hidden kernel and projector are
    \[
    \Kernel{\sigma}
    \equiv
    \ker\bigl(D\SelCoreData{\sigma}\bigr)
    \cap
    \ker\bigl(D\SelCoreShell{\sigma}\bigr),
    \qquad
    \StateComp{\sigma}\equiv \SelCoreProj{\sigma}.
    \]
\end{enumerate}
\end{definition}

\begin{proposition}[A split core reconstructs a selector law]
\label{prop:recovered}
Assume Definitions \ref{def:splitcore}, \ref{def:generator}, and \ref{def:recovered}. Then the recovered data of Definition \ref{def:recovered} satisfy every requirement of Definition \ref{def:generator}.
\end{proposition}

\begin{proof}
The space \(\GenSpace{\sigma}\) in \eqref{eq:recspace} is Euclidean, and the body \(\GenBody{\sigma}\) in \eqref{eq:recbody} is compact and convex with nonempty interior because it is the product of compact convex bodies each with nonempty interior. The readouts \(\GenData{\sigma}\) and \(\GenShell{\sigma}\) are affine by \eqref{eq:recreadouts}. The hidden selector \(\HiddenSel{\sigma}\) from \eqref{eq:reselector} is plainly self-adjoint and idempotent, with kernel \(\SelCoreSpace{\sigma}\) and image \(\GapSpace{\sigma}\).

The exact body splitting \eqref{eq:gensplitting} holds by construction from \eqref{eq:recbody}. The hidden body and hidden operator already satisfy the required positivity conditions from Definition \ref{def:splitcore}. The reconstructed symmetry and time-reversal data preserve all recovered constituents because the split-core symmetries preserve the zero-hidden sector, the hidden sector, the coupling maps, the hidden body, and the hidden operator.

For the shell realization, \(\GenEmbed{\sigma}(\ForSpace{\sigma})\subset\GenBody{\sigma}\cap\ZeroHiddenSlice{\sigma}\) because \(\SelCoreEmbed{\sigma}(\ForSpace{\sigma})\subset\SelCoreBody{\sigma}\). Also
\[
\GenShell{\sigma}\bigl(\GenEmbed{\sigma}(w)\bigr)
=
\SelCoreShell{\sigma}\bigl(\SelCoreEmbed{\sigma}(w)\bigr)
=
0,
\]
so the image lies in the exact selector-zero shell-zero locus. Conversely, every point in
\[
\GenBody{\sigma}\cap\ZeroHiddenSlice{\sigma}\cap
\bigl(\GenShell{\sigma}\bigr)^{-1}(0)
\]
has the form \(u+0\) with \(u\in\SelCoreBody{\sigma}\cap(\SelCoreShell{\sigma})^{-1}(0)\), hence equals \(\SelCoreEmbed{\sigma}(w)+0\) for a unique \(w\in\ForSpace{\sigma}\). Equation \eqref{eq:gencompat} is exactly \eqref{eq:splitcompat}.

Selector-zero solvability \eqref{eq:gensolvability} is just \eqref{eq:splitsolvability} written inside the recovered selector presentation. The inert kernel and projector agree by definition.

For coercivity, let \(w\in\ForSpace{\sigma}\) and \(\delta g=\delta u+\eta\in T_{\GenEmbed{\sigma}(w)}\GenSpace{\sigma}\) satisfy \(D\GenData{\sigma}\,\delta g=0\). Then
\[
D\SelCoreData{\sigma}\,\delta u+\SelHiddenData{\sigma}\eta=0,
\]
and
\[
D\GenShell{\sigma}\,\delta g
=
D\SelCoreShell{\sigma}\,\delta u+\SelHiddenShell{\sigma}\eta,
\qquad
\HiddenSel{\sigma}\delta g=\eta,
\qquad
\StateComp{\sigma}\bigl((I-\HiddenSel{\sigma})\delta g\bigr)
=
\SelCoreProj{\sigma}\delta u.
\]
Applying \eqref{eq:splitcoercive} therefore yields \eqref{eq:gencoercive}. All requirements of Definition \ref{def:generator} are satisfied.
\end{proof}

\section{Collapse and Uniqueness}

\begin{proposition}[The full selector law carries no extra benchmark-facing freedom beyond the split core]
\label{prop:collapse}
Assume Definitions \ref{def:generator}, \ref{def:splitcore}, \ref{def:canonical}, and \ref{def:recovered}. Then for each sector \(\sigma\):
\begin{enumerate}
    \item the zero-hidden state space and body, the zero-hidden observer-data and shell-response readouts, the hidden coupling maps, the hidden body with positive hidden-sector operator, the symmetry and time-reversal structure, the exact zero-response shell realization, fixed-data shell solvability, the inert kernel and projector, and the split coercive control are canonically induced by the full selector law;
    \item once those split-core data are fixed, the larger selector-state space is necessarily the orthogonal product \(\SelCoreSpace{\sigma}\oplus\GapSpace{\sigma}\), the selector-state body is necessarily the product body \eqref{eq:recbody}, the hidden selector is necessarily the product projection \eqref{eq:reselector}, and the full affine selector readouts are necessarily the split zero-hidden readouts plus the hidden coupling maps as in \eqref{eq:recreadouts};
    \item any other selector-law presentation compatible with the same split core differs from the recovered one only by the obvious orthogonal identification preserving the same zero-hidden sector, hidden factor, coupling maps, hidden quadratic form, and benchmark shell data.
\end{enumerate}
\end{proposition}

\begin{proof}
Item (1) is Proposition \ref{prop:induced}. Item (2) is Proposition \ref{prop:recovered}. For item (3), any compatible selector-law presentation must preserve the orthogonal decomposition into the zero-hidden and hidden sectors singled out by the split core, must preserve the product body and the same hidden operator on the hidden factor, and must preserve the same affine zero-hidden readouts and hidden coupling maps. That leaves only the obvious orthogonal identification respecting those data. Therefore the ambient selector-state presentation and the hidden selector projector carry no extra benchmark-facing freedom once the split core is fixed.
\end{proof}

\section{Benchmark Preservation}

\begin{theorem}[The split core preserves the same one-metric benchmark package]
\label{thm:outputs}
Assume Definition \ref{def:generator} or, equivalently by Proposition \ref{prop:recovered}, Definition \ref{def:splitcore}. Then the benchmark output statements of Definition \ref{def:fixed} remain unchanged:
\[
\MetricOut(\mathbf w)=\CompletedMetric,
\qquad
\MatterOut(\mathbf w,\psi)
=
T_{\mu\nu}^{\mathrm{matter}}[\CompletedMetric,\psi]
\]
for every \(\mathbf w\in\PrimShell\) and every matter configuration \(\psi\in\Psi\). Consequently the split-core collapse introduces neither a second operative metric nor an enlarged low-energy matter law.
\end{theorem}

\begin{proof}
The present note changes only the upstream presentation of the selector-law data. It does not alter the primitive shell \(\PrimShell\), the recorded metric map \(\MetricOut\), or the recorded matter map \(\MatterOut\). Those remain the fixed continuation data of Definition \ref{def:fixed}, so \eqref{eq:fixedoutputs} continues to hold exactly.
\end{proof}

\section{Main Theorem}

\begin{theorem}[The remaining burden moves below the full selector-law presentation]
\label{thm:main}
Assume the fixed primitive continuation data of Definition \ref{def:fixed} and, for each sector \(\sigma\in\{\Car,\Cur,\Hom\}\), a zero-hidden selector generator law in the sense of Definition \ref{def:generator}. Then:
\begin{enumerate}
    \item the split zero-hidden selector-shell core of Definition \ref{def:splitcore} is canonically induced by the selector law;
    \item the larger selector-state space, the selector-state body, the self-adjoint hidden selector, the full affine selector readouts, the exact zero-response shell realization, the selector-zero solvability statement, and the selector-law coercivity package are all forced once that split core is fixed;
    \item the one operative metric remains exactly \(\CompletedMetric\), the ordinary matter route remains exactly the standard one through that metric, there is no second operative metric, and the low-energy matter law is not enlarged;
    \item consequently the benchmark-facing burden after the current combined-response-from-selector theorem is discharged in the sharper form now available on the active line: the full zero-hidden selector generator law carries no extra benchmark-facing freedom beyond the smaller split zero-hidden selector-shell core \(\SelectorCore\).
\end{enumerate}
\end{theorem}

\begin{proof}
Item (1) is Proposition \ref{prop:induced}. Item (2) is Proposition \ref{prop:collapse}. Item (3) is Theorem \ref{thm:outputs}. Item (4) is the combined routing consequence of the previous items.
\end{proof}

\begin{corollary}[What remains open after this theorem]
\label{cor:next}
After Theorem \ref{thm:main}, the remaining honest burden is no longer to justify the full zero-hidden selector generator law as such on the active primitive-observer-reduction line. The next honest burden is sharper again: first try to derive or uniquely force the split zero-hidden selector-shell core \(\SelectorCore\) from still more primitive explicit substrate structure in a way that actually changes benchmark-facing status, or else prove a genuinely smaller theorem-level collapse, sharper necessity result, or sharper uniqueness result strictly inside that split core while preserving the same benchmark one-metric package and claim boundary.
\end{corollary}

\begin{proof}
Theorem \ref{thm:main} converts the full selector-law presentation from the remaining benchmark-facing stopping point into a canonical presentation of the smaller split core. What remains open is therefore why that split core itself should hold, rather than becoming the new disciplined stopping point.
\end{proof}

\section{Conclusion}

The positive result of this note is deliberately narrower than a completed final microphysical derivation and stronger than another deeper same-output selector relay. The current selector-law forcing theorem had already shown that the combined response substrate law is a forced reduction of the zero-hidden selector generator law. The present note now spends the legitimate internal-collapse option below that frontier: it shows that the full selector-law presentation itself carries no extra benchmark-facing freedom beyond a smaller split zero-hidden selector-shell core \(\SelectorCore\).

That gain directly addresses the precise constituents that had remained live at selector scope. The ambient selector-state space, the product selector-state presentation, and the self-adjoint hidden selector are no longer independent primitive data. Once the zero-hidden sector, the hidden sector, and their affine coupling readouts are fixed together with the exact shell realization, solvability statement, inert-kernel projector, and split coercive control, the larger selector-law package is recovered canonically.

The benchmark boundary remains fixed throughout. The same primitive shell remains the shell carrying the observer outputs, so the one operative metric is still exactly \(\CompletedMetric\), the ordinary matter route is unchanged, there is no second operative metric, and the low-energy matter law is not enlarged. This is therefore a disciplined upstream derivational gain, not a final completion claim.

The open burden sharpens once more. The note does not yet derive the split selector-shell core from ultimate substrate microphysics. The next benchmark-facing attempt should therefore first try to derive or uniquely force that split core from still more primitive explicit substrate structure in a way that actually changes benchmark-facing status. If that cannot be done cleanly, the admissible fallback is a genuinely smaller collapse, sharper necessity result, or sharper uniqueness result strictly inside that split core. Another full-selector-law restatement, combined-response restatement, ambient-package restatement, trace-core restatement, completion-core restatement, exact-shell affine nucleus restatement, or another same split-core relay would not meet that burden, and no stronger first-principles promotion language is licensed yet.

\end{document}
